\documentclass[12pt]{article}
\usepackage[T1]{fontenc}
\usepackage[utf8]{inputenc}
\usepackage{lmodern}
\usepackage{textcomp}
\usepackage{enumitem}
\usepackage{geometry}
\geometry{margin=1in}
\begin{document}
\begin{center}
Answer Key - Health Sciences - Undergraduate Level Assessment 16 \
\small (Answers are ordered exactly as in the question paper.)
\end{center}

\section*{Multiple Choice}
\begin{enumerate}[label=\arabic*.]
\item \textbf{Answer:} B
\\ \textit{Options:} A) Hygiene; B) Vaccination; C) Antiviral drugs; D) Humanised monoclonal antibodies
\\ \textit{Note:} Vaccination is the primary method for controlling the influenza virus in high-risk populations as it helps to build immunity and reduce the spread of the virus, thus protecting both individuals and the community at large.
\item \textbf{Answer:} D
\\ \textit{Options:} A) masturbation; B) autoeroticism; C) self-eroticism; D) a and b
\\ \textit{Note:} The correct answer is "a and b" because both terms referenced in options a and b refer to the same concept of sexual self-stimulation, commonly known as masturbation.
\item \textbf{Answer:} B
\\ \textit{Options:} A) 11 pairs; B) 23 pairs; C) 32 pairs; D) 46 pairs
\\ \textit{Note:} The correct answer is 23 pairs because human cells typically contain 46 chromosomes organized into 23 pairs, with one chromosome of each pair inherited from each parent.
\item \textbf{Answer:} B
\\ \textit{Options:} A) Alveoli; B) Epiglottis; C) Larynx; D) Uvula
\\ \textit{Note:} The epiglottis is a flap of tissue that covers the trachea during swallowing, preventing food and liquid from entering the lower airway and ensuring that they are directed into the esophagus.
\item \textbf{Answer:} D
\\ \textit{Options:} A) The virion has a dumbbell appearance; B) It is shaped like a bullet from a gun; C) The virus is star shaped; D) The virion is very pleomorphic
\\ \textit{Note:} The correct answer is "the virion is very pleomorphic" because the rabies virus exhibits a variable shape and size, which is characteristic of pleomorphic viruses, distinguishing it from more uniform viral forms.
\end{enumerate}

\section*{True / False}
\begin{enumerate}[label=\arabic*.]
\setcounter{enumi}{5}
\item \textbf{Answer:} True
\\ \textit{Options:} A) True; B) False
\\ \textit{Note:} Individuals with an internal locus of control believe they have the power to influence their circumstances, which often leads to proactive behaviors in seeking solutions and adapting their environments to meet their needs, thereby reducing potential accessibility problems.
\item \textbf{Answer:} False
\\ \textit{Options:} A) True; B) False
\\ \textit{Note:} In vitro fertilization is not necessarily the best choice for women who have experienced repeated early miscarriages because the underlying causes of these miscarriages, such as genetic factors or anatomical issues, may not be addressed by IVF, and alternative treatments or investigations may be more appropriate.
\item \textbf{Answer:} True
\\ \textit{Options:} A) True; B) False
\\ \textit{Note:} The statement is true because the World Health Organization estimates that approximately 350 million people are chronically infected with hepatitis B virus globally, highlighting its prevalence as a major public health concern.
\item \textbf{Answer:} True
\\ \textit{Options:} A) True; B) False
\\ \textit{Note:} By$_{28}$ days after conception, the embryo has undergone significant development and typically measures about 1 centimetre in length as it transitions from the embryonic disc to a more recognizable form.
\item \textbf{Answer:} False
\\ \textit{Options:} A) True; B) False
\\ \textit{Note:} Single men do not represent the largest demographic served by Medicaid; instead, women, children, and low-income families are the primary groups covered by the program.
\end{enumerate}

\section*{Long Form Response}
\begin{enumerate}[label=\arabic*.]
\setcounter{enumi}{10}
\item \textbf{Answer:} Dental plaque is most likely to accumulate on the buccal surfaces of the upper molars and the lingual surfaces of the lower incisors in individuals with poor oral hygiene. The upper molars are particularly susceptible due to their position in the mouth, where food particles and saliva can easily adhere, combined with the fact that they are often less effectively cleaned during brushing. Similarly, the lingual surfaces of the lower incisors are difficult to reach and often neglected, making them prone to plaque buildup. Contributing factors include inadequate brushing technique, infrequent dental care, and dietary habits that favor plaque formation, such as high sugar intake. These areas create environments conducive to bacterial colonization, leading to increased plaque accumulation and potential dental issues.
\\ \textit{Note:} The answer correctly identifies the specific tooth surfaces prone to plaque accumulation due to their anatomical positions and accessibility challenges during oral hygiene practices, emphasizing the impact of poor brushing techniques on plaque retention.
\item \textbf{Answer:} The pupillary light reflex is a consensual reflex that occurs when light is shone into one eye, causing both pupils to constrict simultaneously, regardless of which eye is exposed to the light. This reflex involves the optic nerve transmitting sensory information from the retina to the brain, where it is processed, and the response is executed by the oculomotor nerve controlling the muscles of the iris. The consensual nature of this reflex is significant in assessing neurological function, as it indicates proper communication between the sensory and motor pathways of the nervous system. Abnormalities in the pupillary light reflex can signal potential neurological damage or dysfunction, aiding in the diagnosis of conditions such as brain injury or increased intracranial pressure. Thus, the reflex serves not only as a basic physiological response but also as a critical tool in clinical assessments of neurological health.
\\ \textit{Note:} The answer correctly describes the pupillary light reflex as a consensual response involving both pupils' simultaneous constriction when light is shone in one eye, highlighting the roles of the optic and oculomotor nerves, which is crucial for assessing neurological function.
\end{enumerate}

\vfill
\noindent\textit{End of Answer Key}
\end{document}